\nonumsection{ЗАКЛЮЧЕНИЕ}

Задача эффективного распределения графовых данных между узлами хранения является одной из ключевых проблем при построении масштабируемых графовых баз данных. Рост объёмов данных и увеличение сложности запросов приводят к необходимости минимизировать количество межшардовых взаимодействий, поскольку именно они оказывают существенное влияние на производительность распределённых систем.

В процессе исследования были проанализированы существующие подходы к распределению графовых данных. Рассмотрены методы, основанные на структуре графа, алгоритмы, учитывающие особенности рабочей нагрузки, а также потоковые методы распределения. Проведённый анализ показал, что различные подходы ориентированы на решение разных задач и обладают собственными преимуществами и ограничениями. При этом наиболее перспективными для практического применения являются алгоритмы, способные обеспечивать разумный компромисс между качеством распределения и вычислительными затратами на его построение.

Для решения поставленной задачи был разработан программный модуль распределения вершин в графовой базе данных. Реализованная архитектура обеспечивает возможность использования различных алгоритмов распределения и проведения их сравнительного анализа. Особое внимание было уделено модульности решения, что позволяет расширять функциональность системы и интегрировать новые методы распределения без существенной переработки существующего кода.

Разработанный модуль позволяет выполнять распределение вершин между шардами, анализировать качество полученного распределения и моделировать выполнение запросов в распределённой среде. В качестве основного критерия оценки использовалось количество межшардовых переходов, возникающих при выполнении запросов поиска пути. Данная метрика напрямую связана с объёмом сетевого взаимодействия между узлами и позволяет объективно оценивать эффективность алгоритмов распределения.

Проведённые эксперименты подтвердили работоспособность реализованного решения и показали, что использование специализированных алгоритмов распределения позволяет существенно снизить количество межшардовых переходов по сравнению со случайным размещением данных. Полученные результаты демонстрируют, что учёт структуры графа при распределении данных оказывает значительное влияние на локальность выполнения запросов и общую эффективность системы.

Дополнительное улучшение качества распределения достигается за счёт применения методов локальной оптимизации, позволяющих корректировать первоначальное размещение вершин. Несмотря на увеличение вычислительных затрат, подобный подход позволяет уменьшить количество межшардовых обращений и тем самым повысить эффективность обработки запросов в распределённой среде.

Разработанное решение может использоваться как в исследовательских целях для изучения алгоритмов распределения графовых данных, так и в качестве основы для создания практических систем управления размещением данных в распределённых графовых базах данных. Модульная структура программного обеспечения обеспечивает возможность его дальнейшего развития и адаптации под различные сценарии использования.

Полученные результаты подтверждают эффективность выбранного подхода и показывают перспективность дальнейших исследований в данном направлении. Дальнейшее развитие работы может быть связано с использованием информации о рабочей нагрузке при принятии решений о размещении данных, реализацией механизмов динамического перераспределения графа, исследованием методов контролируемой редундантности, а также разработкой алгоритмов, способных адаптироваться к изменению структуры графа и характера запросов в процессе эксплуатации системы.

Таким образом, разработанный программный модуль позволяет решать задачу распределения вершин в графовой базе данных и создавать распределения, обеспечивающие снижение количества межшардовых переходов при сохранении баланса между узлами хранения. Полученные результаты могут служить основой для дальнейшего совершенствования методов распределённого хранения и обработки графовых данных.
